\documentclass[11pt]{article}
\usepackage[margin=1in]{geometry}
\usepackage[T1]{fontenc}
\usepackage[utf8]{inputenc}
\usepackage{lmodern}
\usepackage{booktabs}
\usepackage{multirow}
\usepackage{graphicx}
\usepackage{amsmath}
\usepackage{amssymb}
\usepackage{microtype}
\usepackage[hidelinks]{hyperref}
\usepackage[numbers,sort&compress]{natbib}

% generated by clef/eval/ft_paper.py -- do not edit
\expandafter\def\csname num:e3:names\endcsname{84}
\expandafter\def\csname num:e3:names_recur\endcsname{3}
\expandafter\def\csname num:e4:consistent\endcsname{100}
\expandafter\def\csname num:e4:lookup_acc\endcsname{99.0}
\expandafter\def\csname num:e4:lookup_cov\endcsname{94}
\expandafter\def\csname num:e4:n_chapters\endcsname{61}
\expandafter\def\csname num:e4:n_entities\endcsname{92}
\expandafter\def\csname num:gain:e3:lenient:acc\endcsname{89.7}
\expandafter\def\csname num:gain:e3:lenient:auc\endcsname{0.955}
\expandafter\def\csname num:gain:e3:released:acc\endcsname{79.2}
\expandafter\def\csname num:gain:e3:released:auc\endcsname{0.900}
\expandafter\def\csname num:gain:e3:strict:acc\endcsname{89.0}
\expandafter\def\csname num:gain:e3:strict:auc\endcsname{0.952}
\expandafter\def\csname num:gain:e4facts:lenient:acc\endcsname{89.4}
\expandafter\def\csname num:gain:e4facts:lenient:f1\endcsname{0.750}
\expandafter\def\csname num:gain:e4facts:released:acc\endcsname{83.6}
\expandafter\def\csname num:gain:e4facts:released:f1\endcsname{0.739}
\expandafter\def\csname num:gain:e4facts:strict:acc\endcsname{88.4}
\expandafter\def\csname num:gain:e4facts:strict:f1\endcsname{0.734}
\expandafter\def\csname num:live:e3:ans_agree_all\endcsname{99.3}
\expandafter\def\csname num:live:e3:big_agree\endcsname{99.0}
\expandafter\def\csname num:live:e3:big_live_acc\endcsname{84.4}
\expandafter\def\csname num:live:e3:big_live_missed\endcsname{5}
\expandafter\def\csname num:live:e3:big_mean\endcsname{534}
\expandafter\def\csname num:live:e3:big_p95\endcsname{619}
\expandafter\def\csname num:live:e3:ci\endcsname{[+1.9, +4.2]}
\expandafter\def\csname num:live:e3:comp_mean\endcsname{224}
\expandafter\def\csname num:live:e3:conf_maxdiff\endcsname{0.0218}
\expandafter\def\csname num:live:e3:d_big\endcsname{+2.9}
\expandafter\def\csname num:live:e3:errors\endcsname{0}
\expandafter\def\csname num:live:e3:esc\endcsname{24.7}
\expandafter\def\csname num:live:e3:esc_agree\endcsname{100.0}
\expandafter\def\csname num:live:e3:lat_err\endcsname{+5.5}
\expandafter\def\csname num:live:e3:live_acc\endcsname{87.3}
\expandafter\def\csname num:live:e3:live_missed\endcsname{5}
\expandafter\def\csname num:live:e3:meas_mean\endcsname{213}
\expandafter\def\csname num:live:e3:meas_p50\endcsname{80}
\expandafter\def\csname num:live:e3:meas_p95\endcsname{624}
\expandafter\def\csname num:live:e3:n\endcsname{864}
\expandafter\def\csname num:live:e3:small_live_acc\endcsname{89.6}
\expandafter\def\csname num:live:e3:small_live_missed\endcsname{10}
\expandafter\def\csname num:live:e3:small_mean\endcsname{99}
\expandafter\def\csname num:live:e3:speedup_meas\endcsname{2.51}
\expandafter\def\csname num:live:e3:speedup_pred\endcsname{2.38}
\expandafter\def\csname num:live:e3:tput_big\endcsname{4.06}
\expandafter\def\csname num:live:e3:tput_cascade\endcsname{10.18}
\expandafter\def\csname num:live:e3:tput_gain\endcsname{2.51}
\expandafter\def\csname num:live:e3:tput_small\endcsname{13.39}
\expandafter\def\csname num:live:e4facts:ans_agree_all\endcsname{99.4}
\expandafter\def\csname num:live:e4facts:big_agree\endcsname{96.9}
\expandafter\def\csname num:live:e4facts:big_live_acc\endcsname{87.8}
\expandafter\def\csname num:live:e4facts:big_mean\endcsname{509}
\expandafter\def\csname num:live:e4facts:big_p95\endcsname{566}
\expandafter\def\csname num:live:e4facts:ci\endcsname{[-0.1, +2.3]}
\expandafter\def\csname num:live:e4facts:comp_mean\endcsname{117}
\expandafter\def\csname num:live:e4facts:conf_maxdiff\endcsname{0.0237}
\expandafter\def\csname num:live:e4facts:d_big\endcsname{+1.1}
\expandafter\def\csname num:live:e4facts:errors\endcsname{0}
\expandafter\def\csname num:live:e4facts:esc\endcsname{4.3}
\expandafter\def\csname num:live:e4facts:esc_agree\endcsname{99.9}
\expandafter\def\csname num:live:e4facts:lat_err\endcsname{+2.9}
\expandafter\def\csname num:live:e4facts:live_acc\endcsname{88.9}
\expandafter\def\csname num:live:e4facts:meas_mean\endcsname{113}
\expandafter\def\csname num:live:e4facts:meas_p50\endcsname{91}
\expandafter\def\csname num:live:e4facts:meas_p95\endcsname{115}
\expandafter\def\csname num:live:e4facts:n\endcsname{1691}
\expandafter\def\csname num:live:e4facts:small_live_acc\endcsname{88.4}
\expandafter\def\csname num:live:e4facts:small_mean\endcsname{95}
\expandafter\def\csname num:live:e4facts:speedup_meas\endcsname{4.49}
\expandafter\def\csname num:live:e4facts:speedup_pred\endcsname{4.36}
\expandafter\def\csname num:live:e4facts:tput_big\endcsname{3.93}
\expandafter\def\csname num:live:e4facts:tput_cascade\endcsname{11.65}
\expandafter\def\csname num:live:e4facts:tput_gain\endcsname{2.96}
\expandafter\def\csname num:live:e4facts:tput_small\endcsname{11.67}
\expandafter\def\csname num:ntrain:e3:lenient\endcsname{468 / 396}
\expandafter\def\csname num:ntrain:e3:strict\endcsname{285 / 238}
\expandafter\def\csname num:ntrain:e4facts:lenient\endcsname{1639 / 1730}
\expandafter\def\csname num:ntrain:e4facts:strict\endcsname{874 / 901}
\expandafter\def\csname num:prereg:sha\endcsname{f4dfe8a1}
\expandafter\def\csname num:prereg:when\endcsname{2026-10-11T21:02:49+0000}
\expandafter\def\csname num:rel:e3:big_live_acc\endcsname{84.4}
\expandafter\def\csname num:rel:e3:big_live_missed\endcsname{5}
\expandafter\def\csname num:rel:e3:big_mean\endcsname{528}
\expandafter\def\csname num:rel:e3:live_acc\endcsname{84.5}
\expandafter\def\csname num:rel:e3:live_missed\endcsname{5}
\expandafter\def\csname num:rel:e3:meas_mean\endcsname{241}
\expandafter\def\csname num:rel:e3:sim_esc\endcsname{33.1}
\expandafter\def\csname num:rel:e3:speedup_meas\endcsname{2.2}
\expandafter\def\csname num:rel:e3:tput_gain\endcsname{2.4}
\expandafter\def\csname num:rel:e4facts:big_live_acc\endcsname{87.8}
\expandafter\def\csname num:rel:e4facts:big_mean\endcsname{502}
\expandafter\def\csname num:rel:e4facts:live_acc\endcsname{87.8}
\expandafter\def\csname num:rel:e4facts:meas_mean\endcsname{278}
\expandafter\def\csname num:rel:e4facts:sim_esc\endcsname{39.9}
\expandafter\def\csname num:rel:e4facts:speedup_meas\endcsname{1.8}
\expandafter\def\csname num:rel:e4facts:tput_gain\endcsname{2.2}
\expandafter\def\csname num:rep:e3:lenient:big_acc\endcsname{84.7}
\expandafter\def\csname num:rep:e3:lenient:big_missed\endcsname{6}
\expandafter\def\csname num:rep:e3:lenient:big_ms\endcsname{528}
\expandafter\def\csname num:rep:e3:lenient:cascade_acc\endcsname{87.5}
\expandafter\def\csname num:rep:e3:lenient:ci\endcsname{[+1.7, +3.9]}
\expandafter\def\csname num:rep:e3:lenient:d_big\endcsname{+2.8}
\expandafter\def\csname num:rep:e3:lenient:esc\endcsname{25}
\expandafter\def\csname num:rep:e3:lenient:mean_ms\endcsname{207}
\expandafter\def\csname num:rep:e3:lenient:missed\endcsname{6}
\expandafter\def\csname num:rep:e3:lenient:small_acc\endcsname{89.7}
\expandafter\def\csname num:rep:e3:lenient:small_missed\endcsname{10}
\expandafter\def\csname num:rep:e3:lenient:speedup\endcsname{2.55}
\expandafter\def\csname num:rep:e3:released:big_acc\endcsname{84.7}
\expandafter\def\csname num:rep:e3:released:big_missed\endcsname{6}
\expandafter\def\csname num:rep:e3:released:big_ms\endcsname{528}
\expandafter\def\csname num:rep:e3:released:cascade_acc\endcsname{84.7}
\expandafter\def\csname num:rep:e3:released:ci\endcsname{[-0.3, +0.3]}
\expandafter\def\csname num:rep:e3:released:d_big\endcsname{+0.0}
\expandafter\def\csname num:rep:e3:released:esc\endcsname{33}
\expandafter\def\csname num:rep:e3:released:mean_ms\endcsname{255}
\expandafter\def\csname num:rep:e3:released:missed\endcsname{6}
\expandafter\def\csname num:rep:e3:released:small_acc\endcsname{79.2}
\expandafter\def\csname num:rep:e3:released:small_missed\endcsname{6}
\expandafter\def\csname num:rep:e3:released:speedup\endcsname{2.07}
\expandafter\def\csname num:rep:e3:strict:big_acc\endcsname{84.7}
\expandafter\def\csname num:rep:e3:strict:big_missed\endcsname{6}
\expandafter\def\csname num:rep:e3:strict:big_ms\endcsname{528}
\expandafter\def\csname num:rep:e3:strict:cascade_acc\endcsname{88.5}
\expandafter\def\csname num:rep:e3:strict:ci\endcsname{[+2.5, +5.3]}
\expandafter\def\csname num:rep:e3:strict:d_big\endcsname{+3.8}
\expandafter\def\csname num:rep:e3:strict:esc\endcsname{16}
\expandafter\def\csname num:rep:e3:strict:mean_ms\endcsname{164}
\expandafter\def\csname num:rep:e3:strict:missed\endcsname{9}
\expandafter\def\csname num:rep:e3:strict:small_acc\endcsname{89.0}
\expandafter\def\csname num:rep:e3:strict:small_missed\endcsname{15}
\expandafter\def\csname num:rep:e3:strict:speedup\endcsname{3.22}
\expandafter\def\csname num:rep:e4facts:lenient:big_acc\endcsname{88.3}
\expandafter\def\csname num:rep:e4facts:lenient:big_ms\endcsname{502}
\expandafter\def\csname num:rep:e4facts:lenient:cascade_acc\endcsname{89.7}
\expandafter\def\csname num:rep:e4facts:lenient:ci\endcsname{[+0.2, +2.7]}
\expandafter\def\csname num:rep:e4facts:lenient:d_big\endcsname{+1.4}
\expandafter\def\csname num:rep:e4facts:lenient:esc\endcsname{1}
\expandafter\def\csname num:rep:e4facts:lenient:mean_ms\endcsname{101}
\expandafter\def\csname num:rep:e4facts:lenient:small_acc\endcsname{89.4}
\expandafter\def\csname num:rep:e4facts:lenient:speedup\endcsname{4.98}
\expandafter\def\csname num:rep:e4facts:released:big_acc\endcsname{88.3}
\expandafter\def\csname num:rep:e4facts:released:big_ms\endcsname{502}
\expandafter\def\csname num:rep:e4facts:released:cascade_acc\endcsname{88.2}
\expandafter\def\csname num:rep:e4facts:released:ci\endcsname{[-0.7, +0.5]}
\expandafter\def\csname num:rep:e4facts:released:d_big\endcsname{-0.1}
\expandafter\def\csname num:rep:e4facts:released:esc\endcsname{40}
\expandafter\def\csname num:rep:e4facts:released:mean_ms\endcsname{295}
\expandafter\def\csname num:rep:e4facts:released:small_acc\endcsname{83.6}
\expandafter\def\csname num:rep:e4facts:released:speedup\endcsname{1.70}
\expandafter\def\csname num:rep:e4facts:strict:big_acc\endcsname{88.3}
\expandafter\def\csname num:rep:e4facts:strict:big_ms\endcsname{502}
\expandafter\def\csname num:rep:e4facts:strict:cascade_acc\endcsname{88.9}
\expandafter\def\csname num:rep:e4facts:strict:ci\endcsname{[-0.6, +1.8]}
\expandafter\def\csname num:rep:e4facts:strict:d_big\endcsname{+0.6}
\expandafter\def\csname num:rep:e4facts:strict:esc\endcsname{4}
\expandafter\def\csname num:rep:e4facts:strict:mean_ms\endcsname{116}
\expandafter\def\csname num:rep:e4facts:strict:small_acc\endcsname{88.4}
\expandafter\def\csname num:rep:e4facts:strict:speedup\endcsname{4.34}
\expandafter\def\csname num:vela:jev\endcsname{31.91}
\expandafter\def\csname num:vela:jev_nox\endcsname{42.55}
\expandafter\def\csname num:vela:n_tasks\endcsname{13}
\expandafter\def\csname num:vela:n_up\endcsname{9}
\expandafter\def\csname num:vf:e3:nox-tuned-lenient:cascade_acc\endcsname{87.5}
\expandafter\def\csname num:vf:e3:nox-tuned-lenient:ci\endcsname{[+1.7, +3.9]}
\expandafter\def\csname num:vf:e3:nox-tuned-lenient:d_big\endcsname{+2.8}
\expandafter\def\csname num:vf:e3:nox-tuned-lenient:esc\endcsname{25}
\expandafter\def\csname num:vf:e3:nox-tuned-lenient:mean_ms\endcsname{207}
\expandafter\def\csname num:vf:e3:nox-tuned-lenient:small_acc\endcsname{89.7}
\expandafter\def\csname num:vf:e3:nox-tuned-lenient:speedup\endcsname{2.55}
\expandafter\def\csname num:vf:e3:nox-tuned-strict:cascade_acc\endcsname{88.5}
\expandafter\def\csname num:vf:e3:nox-tuned-strict:ci\endcsname{[+2.5, +5.3]}
\expandafter\def\csname num:vf:e3:nox-tuned-strict:d_big\endcsname{+3.8}
\expandafter\def\csname num:vf:e3:nox-tuned-strict:esc\endcsname{16}
\expandafter\def\csname num:vf:e3:nox-tuned-strict:mean_ms\endcsname{164}
\expandafter\def\csname num:vf:e3:nox-tuned-strict:small_acc\endcsname{89.0}
\expandafter\def\csname num:vf:e3:nox-tuned-strict:speedup\endcsname{3.22}
\expandafter\def\csname num:vf:e3:nox:cascade_acc\endcsname{84.7}
\expandafter\def\csname num:vf:e3:nox:ci\endcsname{[-0.3, +0.3]}
\expandafter\def\csname num:vf:e3:nox:d_big\endcsname{+0.0}
\expandafter\def\csname num:vf:e3:nox:esc\endcsname{33}
\expandafter\def\csname num:vf:e3:nox:mean_ms\endcsname{255}
\expandafter\def\csname num:vf:e3:nox:small_acc\endcsname{79.2}
\expandafter\def\csname num:vf:e3:nox:small_ms\endcsname{83}
\expandafter\def\csname num:vf:e3:nox:speedup\endcsname{2.07}
\expandafter\def\csname num:vf:e3:vela-0.8b+abs:cascade_acc\endcsname{84.8}
\expandafter\def\csname num:vf:e3:vela-0.8b+abs:ci\endcsname{[-0.5, +0.7]}
\expandafter\def\csname num:vf:e3:vela-0.8b+abs:d_big\endcsname{+0.1}
\expandafter\def\csname num:vf:e3:vela-0.8b+abs:esc\endcsname{59}
\expandafter\def\csname num:vf:e3:vela-0.8b+abs:mean_ms\endcsname{356}
\expandafter\def\csname num:vf:e3:vela-0.8b+abs:small_acc\endcsname{75.6}
\expandafter\def\csname num:vf:e3:vela-0.8b+abs:speedup\endcsname{1.48}
\expandafter\def\csname num:vf:e3:vela-0.8b:cascade_acc\endcsname{84.7}
\expandafter\def\csname num:vf:e3:vela-0.8b:ci\endcsname{[-0.7, +0.6]}
\expandafter\def\csname num:vf:e3:vela-0.8b:d_big\endcsname{+0.0}
\expandafter\def\csname num:vf:e3:vela-0.8b:esc\endcsname{55}
\expandafter\def\csname num:vf:e3:vela-0.8b:mean_ms\endcsname{332}
\expandafter\def\csname num:vf:e3:vela-0.8b:small_acc\endcsname{75.6}
\expandafter\def\csname num:vf:e3:vela-0.8b:small_ms\endcsname{48}
\expandafter\def\csname num:vf:e3:vela-0.8b:speedup\endcsname{1.59}
\expandafter\def\csname num:vf:e3:vela-4b+abs:cascade_acc\endcsname{84.6}
\expandafter\def\csname num:vf:e3:vela-4b+abs:ci\endcsname{[-0.3, +0.0]}
\expandafter\def\csname num:vf:e3:vela-4b+abs:d_big\endcsname{-0.1}
\expandafter\def\csname num:vf:e3:vela-4b+abs:esc\endcsname{60}
\expandafter\def\csname num:vf:e3:vela-4b+abs:mean_ms\endcsname{540}
\expandafter\def\csname num:vf:e3:vela-4b+abs:small_acc\endcsname{76.7}
\expandafter\def\csname num:vf:e3:vela-4b+abs:speedup\endcsname{0.98}
\expandafter\def\csname num:vf:e3:vela-4b-noabs:cascade_acc\endcsname{84.6}
\expandafter\def\csname num:vf:e3:vela-4b-noabs:ci\endcsname{[-0.3, +0.0]}
\expandafter\def\csname num:vf:e3:vela-4b-noabs:d_big\endcsname{-0.1}
\expandafter\def\csname num:vf:e3:vela-4b-noabs:esc\endcsname{52}
\expandafter\def\csname num:vf:e3:vela-4b-noabs:mean_ms\endcsname{516}
\expandafter\def\csname num:vf:e3:vela-4b-noabs:small_acc\endcsname{76.7}
\expandafter\def\csname num:vf:e3:vela-4b-noabs:speedup\endcsname{1.02}
\expandafter\def\csname num:vf:e3:vela-4b:cascade_acc\endcsname{84.6}
\expandafter\def\csname num:vf:e3:vela-4b:ci\endcsname{[-0.3, +0.0]}
\expandafter\def\csname num:vf:e3:vela-4b:d_big\endcsname{-0.1}
\expandafter\def\csname num:vf:e3:vela-4b:esc\endcsname{52}
\expandafter\def\csname num:vf:e3:vela-4b:mean_ms\endcsname{495}
\expandafter\def\csname num:vf:e3:vela-4b:small_acc\endcsname{76.7}
\expandafter\def\csname num:vf:e3:vela-4b:small_ms\endcsname{221}
\expandafter\def\csname num:vf:e3:vela-4b:speedup\endcsname{1.07}
\expandafter\def\csname num:vf:e3:vela-9b+abs:cascade_acc\endcsname{84.5}
\expandafter\def\csname num:vf:e3:vela-9b+abs:ci\endcsname{[-0.7, +0.2]}
\expandafter\def\csname num:vf:e3:vela-9b+abs:d_big\endcsname{-0.2}
\expandafter\def\csname num:vf:e3:vela-9b+abs:esc\endcsname{65}
\expandafter\def\csname num:vf:e3:vela-9b+abs:mean_ms\endcsname{751}
\expandafter\def\csname num:vf:e3:vela-9b+abs:small_acc\endcsname{77.8}
\expandafter\def\csname num:vf:e3:vela-9b+abs:speedup\endcsname{0.70}
\expandafter\def\csname num:vf:e3:vela-9b:cascade_acc\endcsname{84.5}
\expandafter\def\csname num:vf:e3:vela-9b:ci\endcsname{[-0.7, +0.2]}
\expandafter\def\csname num:vf:e3:vela-9b:d_big\endcsname{-0.2}
\expandafter\def\csname num:vf:e3:vela-9b:esc\endcsname{53}
\expandafter\def\csname num:vf:e3:vela-9b:mean_ms\endcsname{689}
\expandafter\def\csname num:vf:e3:vela-9b:small_acc\endcsname{77.8}
\expandafter\def\csname num:vf:e3:vela-9b:small_ms\endcsname{405}
\expandafter\def\csname num:vf:e3:vela-9b:speedup\endcsname{0.77}
\expandafter\def\csname num:vf:e4facts:nox-tuned-lenient:cascade_acc\endcsname{89.7}
\expandafter\def\csname num:vf:e4facts:nox-tuned-lenient:ci\endcsname{[+0.2, +2.7]}
\expandafter\def\csname num:vf:e4facts:nox-tuned-lenient:d_big\endcsname{+1.4}
\expandafter\def\csname num:vf:e4facts:nox-tuned-lenient:esc\endcsname{1}
\expandafter\def\csname num:vf:e4facts:nox-tuned-lenient:mean_ms\endcsname{101}
\expandafter\def\csname num:vf:e4facts:nox-tuned-lenient:small_acc\endcsname{89.4}
\expandafter\def\csname num:vf:e4facts:nox-tuned-lenient:speedup\endcsname{4.98}
\expandafter\def\csname num:vf:e4facts:nox-tuned-strict:cascade_acc\endcsname{88.9}
\expandafter\def\csname num:vf:e4facts:nox-tuned-strict:ci\endcsname{[-0.6, +1.8]}
\expandafter\def\csname num:vf:e4facts:nox-tuned-strict:d_big\endcsname{+0.6}
\expandafter\def\csname num:vf:e4facts:nox-tuned-strict:esc\endcsname{4}
\expandafter\def\csname num:vf:e4facts:nox-tuned-strict:mean_ms\endcsname{116}
\expandafter\def\csname num:vf:e4facts:nox-tuned-strict:small_acc\endcsname{88.4}
\expandafter\def\csname num:vf:e4facts:nox-tuned-strict:speedup\endcsname{4.34}
\expandafter\def\csname num:vf:e4facts:nox:cascade_acc\endcsname{88.2}
\expandafter\def\csname num:vf:e4facts:nox:ci\endcsname{[-0.7, +0.5]}
\expandafter\def\csname num:vf:e4facts:nox:d_big\endcsname{-0.1}
\expandafter\def\csname num:vf:e4facts:nox:esc\endcsname{40}
\expandafter\def\csname num:vf:e4facts:nox:mean_ms\endcsname{295}
\expandafter\def\csname num:vf:e4facts:nox:small_acc\endcsname{83.6}
\expandafter\def\csname num:vf:e4facts:nox:small_ms\endcsname{94}
\expandafter\def\csname num:vf:e4facts:nox:speedup\endcsname{1.70}
\expandafter\def\csname num:vf:e4facts:vela-0.8b+abs:cascade_acc\endcsname{87.9}
\expandafter\def\csname num:vf:e4facts:vela-0.8b+abs:ci\endcsname{[-0.8, -0.1]}
\expandafter\def\csname num:vf:e4facts:vela-0.8b+abs:d_big\endcsname{-0.4}
\expandafter\def\csname num:vf:e4facts:vela-0.8b+abs:esc\endcsname{92}
\expandafter\def\csname num:vf:e4facts:vela-0.8b+abs:mean_ms\endcsname{515}
\expandafter\def\csname num:vf:e4facts:vela-0.8b+abs:small_acc\endcsname{68.2}
\expandafter\def\csname num:vf:e4facts:vela-0.8b+abs:speedup\endcsname{0.97}
\expandafter\def\csname num:vf:e4facts:vela-0.8b:cascade_acc\endcsname{87.9}
\expandafter\def\csname num:vf:e4facts:vela-0.8b:ci\endcsname{[-0.7, -0.1]}
\expandafter\def\csname num:vf:e4facts:vela-0.8b:d_big\endcsname{-0.4}
\expandafter\def\csname num:vf:e4facts:vela-0.8b:esc\endcsname{92}
\expandafter\def\csname num:vf:e4facts:vela-0.8b:mean_ms\endcsname{511}
\expandafter\def\csname num:vf:e4facts:vela-0.8b:small_acc\endcsname{68.2}
\expandafter\def\csname num:vf:e4facts:vela-0.8b:small_ms\endcsname{52}
\expandafter\def\csname num:vf:e4facts:vela-0.8b:speedup\endcsname{0.98}
\expandafter\def\csname num:vf:e4facts:vela-4b+abs:cascade_acc\endcsname{87.5}
\expandafter\def\csname num:vf:e4facts:vela-4b+abs:ci\endcsname{[-1.6, +0.0]}
\expandafter\def\csname num:vf:e4facts:vela-4b+abs:d_big\endcsname{-0.8}
\expandafter\def\csname num:vf:e4facts:vela-4b+abs:esc\endcsname{30}
\expandafter\def\csname num:vf:e4facts:vela-4b+abs:mean_ms\endcsname{382}
\expandafter\def\csname num:vf:e4facts:vela-4b+abs:small_acc\endcsname{83.6}
\expandafter\def\csname num:vf:e4facts:vela-4b+abs:speedup\endcsname{1.31}
\expandafter\def\csname num:vf:e4facts:vela-4b-noabs:cascade_acc\endcsname{87.4}
\expandafter\def\csname num:vf:e4facts:vela-4b-noabs:ci\endcsname{[-1.7, -0.1]}
\expandafter\def\csname num:vf:e4facts:vela-4b-noabs:d_big\endcsname{-0.9}
\expandafter\def\csname num:vf:e4facts:vela-4b-noabs:esc\endcsname{26}
\expandafter\def\csname num:vf:e4facts:vela-4b-noabs:mean_ms\endcsname{360}
\expandafter\def\csname num:vf:e4facts:vela-4b-noabs:small_acc\endcsname{83.5}
\expandafter\def\csname num:vf:e4facts:vela-4b-noabs:speedup\endcsname{1.39}
\expandafter\def\csname num:vf:e4facts:vela-4b:cascade_acc\endcsname{87.4}
\expandafter\def\csname num:vf:e4facts:vela-4b:ci\endcsname{[-1.7, -0.1]}
\expandafter\def\csname num:vf:e4facts:vela-4b:d_big\endcsname{-0.9}
\expandafter\def\csname num:vf:e4facts:vela-4b:esc\endcsname{26}
\expandafter\def\csname num:vf:e4facts:vela-4b:mean_ms\endcsname{364}
\expandafter\def\csname num:vf:e4facts:vela-4b:small_acc\endcsname{83.6}
\expandafter\def\csname num:vf:e4facts:vela-4b:small_ms\endcsname{232}
\expandafter\def\csname num:vf:e4facts:vela-4b:speedup\endcsname{1.38}
\expandafter\def\csname num:vf:e4facts:vela-9b+abs:cascade_acc\endcsname{88.2}
\expandafter\def\csname num:vf:e4facts:vela-9b+abs:ci\endcsname{[-0.4, +0.2]}
\expandafter\def\csname num:vf:e4facts:vela-9b+abs:d_big\endcsname{-0.1}
\expandafter\def\csname num:vf:e4facts:vela-9b+abs:esc\endcsname{78}
\expandafter\def\csname num:vf:e4facts:vela-9b+abs:mean_ms\endcsname{816}
\expandafter\def\csname num:vf:e4facts:vela-9b+abs:small_acc\endcsname{82.5}
\expandafter\def\csname num:vf:e4facts:vela-9b+abs:speedup\endcsname{0.62}
\expandafter\def\csname num:vf:e4facts:vela-9b:cascade_acc\endcsname{88.1}
\expandafter\def\csname num:vf:e4facts:vela-9b:ci\endcsname{[-0.6, +0.2]}
\expandafter\def\csname num:vf:e4facts:vela-9b:d_big\endcsname{-0.2}
\expandafter\def\csname num:vf:e4facts:vela-9b:esc\endcsname{72}
\expandafter\def\csname num:vf:e4facts:vela-9b:mean_ms\endcsname{784}
\expandafter\def\csname num:vf:e4facts:vela-9b:small_acc\endcsname{82.5}
\expandafter\def\csname num:vf:e4facts:vela-9b:small_ms\endcsname{424}
\expandafter\def\csname num:vf:e4facts:vela-9b:speedup\endcsname{0.64}

\newcommand{\N}[1]{\csname num:#1\endcsname}
\newcommand{\TODO}[1]{\textbf{[#1]}}
\emergencystretch=2em

\title{Specialize the First Stage:\\
Fine-Tuning a Cascade's Decision Model on One Game Master's Rulings,\\
Predicted by Replay and Checked Live}

\author{Kostadis Roussos\\
Independent researcher\\
\texttt{kostadis@gmail.com}}

\date{October 2026}

\begin{document}
\maketitle

\begin{abstract}
A cascade's savings are set by its first stage: every item the cheap model cannot settle is paid
for twice. Two earlier studies on a tabletop role-playing campaign pipeline evaluated off-the-shelf
non-generative decision models against a generative model~\cite{companion}, and designed a
Decision-2.0-Nox-4B $\rightarrow$ \texttt{qwen3.8-flash-next} cascade by replaying their frozen
outputs, then checked it live~\cite{cascade}. The first study also fine-tuned Nox on the game
master's (GM's) rulings but never put the tuned model in the cascade. Here we do, asking whether a
domain-specialized first stage removes more of the second. Because the frozen dataset caches Nox's
backbone features, tuning its decision head and replaying the resulting cascade takes about a minute
of CPU time per variant, cross-fitted so that no item's label reaches its own head or its own gate.
On triage, the tuned head lowered escalation from \N{rep:e3:released:esc}\% to
\N{rep:e3:lenient:esc}\% and raised the cascade's predicted single-stream speed-up from
\N{rep:e3:released:speedup}$\times$ to \N{rep:e3:lenient:speedup}$\times$, while still missing no more
real names than the generative model. On fact typing the replay first showed the second stage
almost vanishing, then exposed why: every fact's label is its entity's label and entities recur, so
a plain lookup of the GM's earlier rulings scores \N{e4:lookup_acc}\%, above every model. Held out
by entity, the tuned Nox alone tied the generative model on entities it had never seen
(\N{rep:e4facts:strict:small_acc}\% vs \N{rep:e4facts:strict:big_acc}\%). Holding out by token on
triage broke the missed-name constraint. We then served the tuned heads on the real two-box
system, with gates and predictions pre-registered by hash. Live, the triage cascade scored
\N{live:e3:d_big} points above qwen (CI \N{live:e3:ci}; predicted \N{rep:e3:lenient:d_big}), escalated
\N{live:e3:esc}\% of tokens, missed exactly as many names as qwen, and was
\N{live:e3:speedup_meas}$\times$ faster single-stream (predicted \N{live:e3:speedup_pred}$\times$); the
entity-disjoint fact-typing cascade escalated \N{live:e4facts:esc}\% and was
\N{live:e4facts:speedup_meas}$\times$ faster (predicted \N{live:e4facts:speedup_pred}$\times$) at
\N{live:e4facts:d_big} points (CI \N{live:e4facts:ci}). The replay routed
\N{live:e3:esc_agree}\% and \N{live:e4facts:esc_agree}\% of items as the live system did. A
specialized first stage removes most of the second on fact typing and a quarter of it on triage,
where the second stage stays because the tuned model alone drops twice as many real names. Set
beside Vela 2.0, a full fine-tune of the same Nox for safety and routing, the head-only gain on our
tasks (+5.5 AUC points on triage) is of the same order as Vela's mid-range gains on its own, at a
small fraction of the effort; Vela itself, tuned for another domain, is a worse first stage here
than the Nox it came from.
\end{abstract}

\section{Introduction}

A reader of the cascade study~\cite{cascade} asked the natural follow-up: if the small model's
accuracy gap is what forces escalation, why accept its baseline accuracy? For a bounded
classification task one can trade general-purpose performance for domain specialization by
fine-tuning, and a specialized first stage might close enough of the gap to remove the second stage
for more of the work. The question has a practical edge in this pipeline. The decision model answers
in about 80--150~ms on a DGX Spark and the generative model in about half a second, so each
percentage point of escalation removed is latency and queue capacity returned.

Both ingredients already existed. The companion study's E6~\cite{companion} fine-tuned Nox's head
(and a LoRA) on the GM's rulings and reported that head-only tuning \emph{recalibrates} the readout:
triage AUC rose from 0.900 to 0.963. But it scored the tuned model as a standalone classifier. The
cascade study's replay~\cite{cascade} showed that a cascade's accuracy, escalation and latency are a
pure function of saved per-item outputs when the first stage is deterministic, and that the
prediction held live. Here we combine them, with three aims:
\begin{enumerate}
  \item measure how much a tuned first stage shrinks the second, at the cascade's own objective
        (the generative model's accuracy, with real names never dropped more often than it drops them);
  \item find where the gain comes from, and in particular whether it is specialization or
        memorization of rulings the GM has already made;
  \item check the replay's prediction for a \emph{trained} first stage on the real system, as the
        cascade study did for an off-the-shelf one.
\end{enumerate}

\paragraph{Findings.}
\begin{itemize}
  \item \textbf{Triage.} Tuning removes a quarter of the escalations (\N{rep:e3:released:esc}\% to
        \N{live:e3:esc}\%) at the missed-name constraint, and the cascade ends up
        \N{live:e3:d_big} points \emph{more} accurate than qwen alone, live. The second stage stays:
        tuned Nox alone is more accurate still (\N{live:e3:small_live_acc}\%) but drops
        \N{live:e3:small_live_missed} real names to qwen's \N{live:e3:big_live_missed}.
  \item \textbf{Fact typing.} Most of the apparent gain is memorization: the label is an
        entity-level ruling, and a lookup of earlier rulings scores \N{e4:lookup_acc}\%. Held out by
        entity, the gain that remains is real but small. Tuned Nox alone ties qwen
        (\N{live:e4facts:small_live_acc}\% vs \N{live:e4facts:big_live_acc}\% live), so the cascade
        escalates only \N{live:e4facts:esc}\% and runs \N{live:e4facts:speedup_meas}$\times$ faster
        than qwen.
  \item \textbf{The replay held for a trained first stage.} Escalation, accuracy difference and
        single-stream latency were all predicted before any live request, and all landed within
        about 5\% (latency) or half a point (accuracy).
  \item \textbf{Next to Vela 2.0.} Head-only tuning on one GM's rulings gained about as much on our
        tasks as Vela's full fine-tune gained on its mid-range tasks (+5.5 AUC points on triage), but
        on fact typing the gain is accuracy on the common types, not macro-F1. Served as a first
        stage on our tasks, no Vela member improved on released Nox.
  \item \textbf{One token-disjoint variant failed} the missed-name constraint on the replay and
        never went live; the replay is where that should be found.
\end{itemize}

\section{Background}

\paragraph{The pipeline and its frozen dataset.} The companion study~\cite{companion} evaluated
decision models, which read a \texttt{state} once and return a probability for every option of a
typed question through a small head without generating text, on judgment steps of a pipeline that
turns recorded play sessions into campaign records. Its labels are one GM's past rulings, frozen
into a versioned, read-only dataset with every model-facing input pre-rendered. We use two of its
tasks, unchanged:
\begin{description}
  \item[E3 triage] (864 tokens, \N{e3:names} real names) is an unknown capitalised token a campaign
        name, or filler, a rules term or real-world chatter? As in the cascade study we score it as
        binary accuracy and separately count \emph{missed names}, because a dropped name is a silent
        deletion nothing downstream recovers.
  \item[E4 fact typing] (1{,}691 facts, \N{e4:n_entities} entities, \N{e4:n_chapters} chapters) file
        an extracted fact under one of eight types.
\end{description}
The other two tasks of the cascade study are not candidates. On E4 entity typing Nox alone already
matched the generative model, and on E2 spelling Nox trailed it by 56 points, a gap that looks
structural rather than one of calibration.

\paragraph{The cascade and its replay.} Nox answers first; if its gate confidence (top probability,
or $\max(p, 1-p)$ for the binary triage question) is below a threshold, the item goes to
\texttt{qwen3.8-flash-next} with reasoning off. The \emph{parity gate} is the lowest threshold whose
cascade accuracy on a training set is at least qwen's, and on E3 misses no more names than qwen.
Thresholds are fitted 2-fold by group (sessions for E3, chapters for E4) and each half is scored with
the other half's threshold. The replay composes per-item answers, confidences and single-stream
service times; the cascade study found it routed items exactly as the live system did and predicted
the cascade-minus-qwen accuracy difference within 0.1 points~\cite{cascade}.

\section{Method}

\subsection{Tuning the head on cached features}
The released Transformers wrapper refuses \texttt{train()}, but the runtime's backbone and head are
ordinary PyTorch modules. As in E6, we ran Nox's frozen backbone once over every item through the
runtime's own loader and renderer and cached the option and query hidden rows the head reads; the
released head on those features reproduces the served probabilities (maximum absolute difference
0.0097 on E3 and 0.0 on E4 in this run). Training is head-only, on CPU: valid-option cross-entropy
($-\log\sum_{k\in\text{valid}}p_k$) at the server's temperature, an L2 pull toward the released head
(L2-SP~\cite{l2sp}, weight 0.1), AdamW at $10^{-4}$, batch 32, 20 epochs, and weight 4 on real-name tokens in
triage. These settings are E6's and were fixed before this study; E6 itself chose them from four
candidates on a cross-validation over the same items (different folds), so they carry a little
optimism. Training is seeded and deterministic: re-training reproduced every probability exactly.
E4 training also uses the companion study's 839-fact control set, whose 54 entities do not occur in
the 1{,}691 test facts.

\subsection{Cross-fitting the head and the gate together}
A tuned first stage adds a second place where labels can leak: the head as well as the threshold.
We nest both inside the replay's existing 2-fold split. For each half $h$:
\begin{itemize}
  \item the \emph{test head} for $h$ is trained on the other half only, and gives the first-stage
        answers and confidences for $h$;
  \item the \emph{gate} for $h$ is the parity gate fitted on the other half, using that half's
        \emph{inner out-of-fold} predictions (3-fold by group within it), so that, like the test
        predictions, they come from heads that never saw the item's label.
\end{itemize}
No item's label therefore influences either its own answer or its own threshold.

\subsection{Two hold-out regimes}
Holding out by session or chapter is the deployment question (\emph{train on the GM's past
sessions, run on the next}), but entities and tokens recur across sessions. We therefore also run a
\emph{disjoint} variant that drops from each training set every item whose entity (E4) or token
(E3, case-folded) occurs in the test half. It answers a different question: does the gain transfer
to things the GM has never ruled on? On E4 it keeps \N{ntrain:e4facts:strict} training items per
half (mostly the control set) against \N{ntrain:e4facts:lenient}; on E3,
\N{ntrain:e3:strict} against \N{ntrain:e3:lenient}.

\subsection{Serving a tuned head, and the live check}
The runtime verifies its model package against a manifest of per-file hashes and a built-in model
identity, and runs a golden-answer check before reporting ready, so a tuned head file cannot simply
be dropped into the package. Our server wrapper lets the stock runtime load, verify and golden-check
the \emph{released} package, then copies the tuned head's weights into the loaded head module in
place and logs the file's hash. Backbone, renderer, scheduler and HTTP API are the stock runtime's.
Each half needs its own head, so the live run serves four heads in turn on spark2, re-probing 20
items per head against the replay's probabilities before measuring. The second stage is qwen on
spark1 in its deployed configuration (MTP-2 speculative decoding, prefix caching on), as in the
cascade study.

We took one variant per task live, chosen on the replay before any live request: tuned on E3
(the only tuned E3 variant that held the missed-name constraint) and disjoint on E4 (the variant
whose gain is not memorization). Every gate, every head's hash and the replay's predictions were
written to a pre-registration record (SHA-256 \texttt{\N{prereg:sha}\ldots}, \N{prereg:when}) before
the first live request. As in the cascade study, we ran tuned Nox alone, the cascade, and qwen alone
on every item in the same session, single-stream ($c=1$) and with 8 concurrent clients ($c=8$).

\section{Replay results}

\begin{table}[t]
\centering
\small
\setlength{\tabcolsep}{4pt}
\resizebox{\textwidth}{!}{\begin{tabular}{llrrrrrr}
\toprule
Task & First stage & Nox alone & Cascade & $\Delta$ vs qwen (pt) [95\% CI] & Esc.\ (\%) & Missed & Speed-up \\
\midrule
E3 triage & released & 79.2 & 84.7 & +0.0 [-0.3, +0.3] & 33 & 6 / 6 & 2.07$\times$ \\
 & tuned & 89.7 & 87.5 & +2.8 [+1.7, +3.9] & 25 & 6 / 6 & 2.55$\times$ \\
 & tuned, disjoint & 89.0 & 88.5 & +3.8 [+2.5, +5.3] & 16 & 9 / 6 & 3.22$\times$ \\
\midrule
E4 fact typing & released & 83.6 & 88.2 & -0.1 [-0.7, +0.5] & 40 & -- & 1.70$\times$ \\
 & tuned & 89.4 & 89.7 & +1.4 [+0.2, +2.7] & 1 & -- & 4.98$\times$ \\
 & tuned, disjoint & 88.4 & 88.9 & +0.6 [-0.6, +1.8] & 4 & -- & 4.34$\times$ \\
E4 fact typing & entity lookup & \multicolumn{2}{r}{99.0 (coverage 94\%)} & & & & \\
\bottomrule
\end{tabular}
}
\caption{Replay of the Nox $\rightarrow$ qwen cascade with the released and the tuned head, held out
as in the text (head and gate cross-fitted). qwen alone scores \N{rep:e3:released:big_acc}\% on E3
and \N{rep:e4facts:released:big_acc}\% on E4. ``Missed'' is real names dropped by the cascade / by
qwen alone. Speed-up is against qwen alone, single-stream, from service times measured live in the
cascade study. ``Disjoint'' also removes from training every item sharing an entity (E4) or token
(E3) with the test half. The lookup row files each fact under the type the GM gave its entity in the
other half, falling back to released Nox when the entity is new.}
\label{tab:replay}
\end{table}

\paragraph{Triage: a smaller second stage.} Tuned on the other half's sessions, Nox alone rose from
\N{rep:e3:released:small_acc}\% to \N{rep:e3:lenient:small_acc}\% binary accuracy, above qwen's
\N{rep:e3:released:big_acc}\% (Table~\ref{tab:replay}). That does not make it a replacement: alone it
drops \N{rep:e3:lenient:small_missed} real names where qwen drops \N{rep:e3:lenient:big_missed}, and
the binary score rewards confident ``not a name'' answers, of which there are many. The parity gate
handles exactly this. The tuned cascade escalates \N{rep:e3:lenient:esc}\% of tokens instead of
\N{rep:e3:released:esc}\%, misses \N{rep:e3:lenient:missed} names (qwen
\N{rep:e3:lenient:big_missed}), scores \N{rep:e3:lenient:d_big} points above qwen (CI
\N{rep:e3:lenient:ci}), and is predicted \N{rep:e3:lenient:speedup}$\times$ faster than qwen
single-stream, against \N{rep:e3:released:speedup}$\times$ for the released head.

This is not name memorization. Only \N{e3:names_recur} of the \N{e3:names} real names have a token
that also occurs in the other half. What recurs is the table's chatter: filler words, rules terms
and real-world references come back every session, and learning how this GM labels them is
legitimate domain knowledge.

\paragraph{Triage, token-disjoint: the constraint breaks.} Removing every training token that also
occurs in the test half cuts the training set to \N{ntrain:e3:strict} items per half. The replay
then escalates only \N{rep:e3:strict:esc}\% and scores higher still
(\N{rep:e3:strict:d_big} points), but the cascade misses \N{rep:e3:strict:missed} names against
qwen's \N{rep:e3:strict:big_missed}. Its gate was fitted to miss no more names than qwen on the
other half and did not transfer. With \N{e3:names} names in all, three names is the difference, but
it is the error this pipeline cannot afford, so this variant fails the objective.

\paragraph{Fact typing: the second stage nearly vanishes, for the wrong reason.} Tuned on the other
half's chapters, Nox alone reaches \N{rep:e4facts:lenient:small_acc}\%, the parity gate escalates
\N{rep:e4facts:lenient:esc}\% of facts, and the cascade is predicted
\N{rep:e4facts:lenient:speedup}$\times$ faster than qwen at \N{rep:e4facts:lenient:d_big} points
above it. The frozen dataset shows why. In E4, \N{e4:consistent}\% of facts carry their entity's
majority label: the label is an entity-level ruling copied onto each fact. The same entities recur
across chapters, so a model trained on other chapters has seen most test entities ruled on already.
A lookup that files each fact under the type the GM gave its entity in the other half covers
\N{e4:lookup_cov}\% of facts and scores \N{e4:lookup_acc}\%, above qwen, tuned Nox and every model in
either earlier study. Much of the tuned head's gain is that lookup learned the hard way, and done
worse.

\paragraph{Fact typing, entity-disjoint: specialization that transfers.} Trained only on entities
absent from the test half (largely the control set), tuned Nox alone scores
\N{rep:e4facts:strict:small_acc}\%, against \N{rep:e4facts:released:small_acc}\% for the released
head and \N{rep:e4facts:strict:big_acc}\% for qwen. The gain here cannot come from remembering
rulings; it is the readout recalibration E6 described, learned from other entities. The parity gate
escalates \N{rep:e4facts:strict:esc}\%, and the predicted speed-up is
\N{rep:e4facts:strict:speedup}$\times$. Its gate rests on few out-of-fold items (the disjoint
regime leaves almost no repair-set training data inside each half), and one half's threshold is 0,
so this point is best read as ``tuned Nox alone ties qwen on unseen entities'' rather than as a
finely tuned cascade.

\section{Live results}

\begin{table}[t]
\centering
\small
\setlength{\tabcolsep}{4pt}
\resizebox{\textwidth}{!}{\begin{tabular}{llrrrrrrr}
\toprule
Task & First stage & Cascade & qwen & $\Delta$ (pt) & Esc. & Missed & Speed-up & Throughput \\
 & & acc.\ (\%) & acc.\ (\%) & [95\% CI] & (\%) & (casc./qwen) & $c{=}1$ & gain $c{=}8$ \\
\midrule
E3 triage & released$^\dagger$ & 84.5 & 84.4 & & 33 & 5 / 5 & 2.2$\times$ & 2.4$\times$ \\
 & tuned & 87.3 & 84.4 & +2.9 [+1.9, +4.2] & 25 & 5 / 5 & 2.51$\times$ & 2.51$\times$ \\
\midrule
E4 fact typing & released$^\dagger$ & 87.8 & 87.8 & & 40 & -- & 1.8$\times$ & 2.2$\times$ \\
 & tuned & 88.9 & 87.8 & +1.1 [-0.1, +2.3] & 4 & -- & 4.49$\times$ & 2.96$\times$ \\
\bottomrule
\end{tabular}
}
\caption{Live Nox $\rightarrow$ qwen cascade on the real system (\N{live:e3:n} triage tokens,
\N{live:e4facts:n} facts, \N{live:e3:errors} errors). Tuned rows: this study, both halves pooled,
each half served with its own head and gate. $^\dagger$Released rows: the cascade study's live run of
the same system with the released head~\cite{cascade}. qwen accuracy and speed-ups are against qwen
alone measured in the same session. Throughput gain is items per second with 8 concurrent clients,
cascade over qwen alone.}
\label{tab:live}
\end{table}

\begin{table}[t]
\centering
\small
\setlength{\tabcolsep}{4pt}
\resizebox{\textwidth}{!}{\begin{tabular}{lrrrrrrrr}
\toprule
Task & Esc.\ pred. & Esc.\ agree & max $|\Delta p|$ & Answers agree & qwen self-agree & $\Delta$ pred. & $\Delta$ meas. & Latency err. \\
\midrule
E3 triage & 24.7\% & 100.0\% & 0.0218 & 99.3\% & 99.0\% & +2.8 & +2.9 & +5\% \\
E4 fact typing & 4.2\% & 99.9\% & 0.0237 & 99.4\% & 96.9\% & +0.6 & +1.1 & +3\% \\
\bottomrule
\end{tabular}
}
\caption{Pre-registered replay predictions against the live run. ``Esc.\ agree'': items the live
cascade routed as the frozen gate on replay probabilities did. $\max|\Delta p|$: largest difference
between the served and the replay's first-stage probability over all items. ``Answers agree'': live
cascade answers equal to the replay's. ``qwen self-agree'': qwen's live answers equal to its frozen
ones, the floor on how well any prediction involving qwen can match. Latency error: composed
single-stream mean against measured.}
\label{tab:check}
\end{table}

\paragraph{Accuracy and routing.} Both cascades did what the replay said
(Tables~\ref{tab:live} and~\ref{tab:check}). On triage the cascade scored \N{live:e3:live_acc}\%
against qwen's \N{live:e3:big_live_acc}\%, \N{live:e3:d_big} points (CI \N{live:e3:ci}) against a
predicted \N{rep:e3:lenient:d_big}, escalated \N{live:e3:esc}\% of tokens, and missed
\N{live:e3:live_missed} real names to qwen's \N{live:e3:big_live_missed}. Tuned Nox alone scored
\N{live:e3:small_live_acc}\% and missed \N{live:e3:small_live_missed}, as on the replay. On fact typing
the cascade scored \N{live:e4facts:live_acc}\% against \N{live:e4facts:big_live_acc}\%
(\N{live:e4facts:d_big} points, CI \N{live:e4facts:ci}, predicted \N{rep:e4facts:strict:d_big}), with
Nox alone at \N{live:e4facts:small_live_acc}\%. The live system escalated exactly the items the frozen
gates selected on E3, and all but one of 1{,}691 on E4. Served and replay probabilities differed by
at most \N{live:e4facts:conf_maxdiff}, slightly above the per-head probe's 0.02 tolerance on a few
items, and the one E4 item near its threshold flipped. The remaining answer disagreement is qwen's:
it reproduced its own frozen answers on \N{live:e3:big_agree}\% and \N{live:e4facts:big_agree}\% of
items.

\paragraph{Latency and throughput.} Single-stream, the triage cascade averaged
\N{live:e3:meas_mean}~ms per token against qwen's \N{live:e3:big_mean}~ms,
\N{live:e3:speedup_meas}$\times$ faster (predicted \N{live:e3:speedup_pred}$\times$, mean
\N{live:e3:lat_err}\% off); the released head's cascade had managed \N{rel:e3:speedup_meas}$\times$.
Fact typing averaged \N{live:e4facts:meas_mean}~ms against \N{live:e4facts:big_mean}~ms,
\N{live:e4facts:speedup_meas}$\times$ (predicted \N{live:e4facts:speedup_pred}$\times$,
\N{live:e4facts:lat_err}\% off), against \N{rel:e4facts:speedup_meas}$\times$ for the released head.
With 8 concurrent clients the cascades sustained \N{live:e3:tput_gain}$\times$ and
\N{live:e4facts:tput_gain}$\times$ qwen's throughput. On fact typing that is all the first stage can
give: the cascade's throughput (\N{live:e4facts:tput_cascade} items/s) equals tuned Nox alone's
(\N{live:e4facts:tput_small}), so the remaining bottleneck is the decision model, not qwen.

\paragraph{A contaminated first pass.} The pre-registered plan measured everything in one session.
In the first pass, another client of the author's began sending long generation requests to spark1's
qwen partway through: it overlapped the qwen-alone runs on both tasks and the second E4 half's
cascade, inflating qwen's single-stream mean to 881 and 808~ms (against
\N{rel:e3:big_mean} and \N{rel:e4facts:big_mean}~ms in the cascade study) and the triage latency
error to +45\%. Accuracy and routing were essentially unchanged. We reran exactly those legs, with the
same heads and gates, while sampling spark1's running-request count every 2~s; no leg ever had more
requests in flight than its own concurrency, and qwen's means returned to \N{live:e3:big_mean} and
\N{live:e4facts:big_mean}~ms. Tables~\ref{tab:live} and~\ref{tab:check} use the rerun; the
contaminated outputs are kept with the results. The episode is the cascade study's caveat in
miniature: the replay composes service times from a quiet box, and a shared second stage breaks that
assumption without telling anyone.

\section{The tuning outcome next to Vela 2.0's}
\label{sec:vela}

Vela 2.0~\cite{vela2} is a published example of the trade the question proposed: vLLM Semantic Router
and KR Labs fine-tuned the same Decision-2.0-Nox-4B for safety, prompt-attack, PII, hallucination and
request-routing decisions. Its tuning is a full fine-tune of 4{,}000 steps on a large public and
synthetic mixture (plus two span-head stages); ours is the decision head alone, on a few hundred to a
couple of thousand of one GM's rulings, in about a minute of CPU time. Table~\ref{tab:gains} sets the
two outcomes side by side, each on its own domain and in the metrics Vela reports: AUC for yes/no
decisions, macro-F1 for multi-class ones.

\begin{table}[t]
\centering
\small
\setlength{\tabcolsep}{4pt}
\resizebox{\textwidth}{!}{\begin{tabular}{lllrrl}
\toprule
Model & Task (held out) & Metric & Before & After & $\Delta$ (pt) \\
\midrule
Nox-4B, tuned head & E3 triage, by session & AUC & 0.900 & 0.955 & +5.5 \\
Nox-4B, tuned head & E3 triage, by session & accuracy (\%) & 79.2 & 89.7 & +10.5 \\
Nox-4B, tuned head & E4 fact typing, by entity & macro-F1 & 0.739 & 0.734 & -0.4 \\
Nox-4B, tuned head & E4 fact typing, by entity & accuracy (\%) & 83.6 & 88.4 & +4.7 \\
Nox-4B, tuned head & E4 fact typing, by chapter & macro-F1 & 0.739 & 0.750 & +1.1 \\
\midrule
Vela 2.0 4B$^\ddagger$ & Prompt attacks, unseen families & AUC & 0.792 & 0.985 & +19.4 [+15.1, +23.7] \\
Vela 2.0 4B$^\ddagger$ & Hate speech, Multilingual HateCheck & AUC & 0.646 & 0.769 & +12.3 [+11.7, +12.9] \\
Vela 2.0 4B$^\ddagger$ & Request harm, CultureGuard & AUC & 0.863 & 0.898 & +3.5 [+3.2, +3.8] \\
Vela 2.0 4B$^\ddagger$ & Request harm, RTP-LX & AUC & 0.761 & 0.774 & +1.3 [+0.8, +1.8] \\
Vela 2.0 4B$^\ddagger$ & Fact-check needed & macro-F1 & 0.501 & 1.000 & +49.9 [+46.1, +53.5] \\
Vela 2.0 4B$^\ddagger$ & Feedback & macro-F1 & 0.425 & 0.920 & +49.5 [+47.0, +52.2] \\
Vela 2.0 4B$^\ddagger$ & Modality & macro-F1 & 0.889 & 1.000 & +11.1 [+9.4, +12.7] \\
Vela 2.0 4B$^\ddagger$ & Domain & macro-F1 & 0.831 & 0.820 & $-$1.2 [$-$2.8, +0.5] \\
Vela 2.0 4B$^\ddagger$ & PII, short texts & F1 & 0.976 & 0.987 & +1.1 [+0.2, +2.1] \\
Vela 2.0 4B$^\ddagger$ & PII, 8K-token documents & F1 & 0.908 & 0.905 & $-$0.3 [$-$3.6, +3.3] \\
Vela 2.0 4B$^\ddagger$ & Hallucination spans & F1 & 0.875 & 0.874 & $-$0.1 [$-$0.7, +0.6] \\
Vela 2.0 4B$^\ddagger$ & RAGTruth & F1 & 0.743 & 0.770 & +2.7 [+0.4, +5.1] \\
Vela 2.0 4B$^\ddagger$ & Toxic spans & char-F1 & 0.451 & 0.379 & $-$7.2 [$-$8.0, $-$6.5] \\
\bottomrule
\end{tabular}
}
\caption{What tuning bought on the tuned-for domain. Tuned Nox: this study, Nox alone, from the
cross-fitted held-out probabilities (no item scored by a head that saw its label); ``before'' is the
released Nox. $^\ddagger$Vela 2.0 4B: as published in its evaluation card~\cite{vela2} (revision
\texttt{d692394}), not re-measured; ``before'' is the per-task reference model the card compares
against (an earlier specialist, a guard model or a span detector), not Nox, whose scores on these
tasks are not published. $\Delta$ in points ($\times 100$ for AUC and F1); Vela's intervals are the
card's paired 95\% CIs.}
\label{tab:gains}
\end{table}

\paragraph{Same order of gain where the comparison is fair.} On triage, tuning raised Nox's AUC from
\N{gain:e3:released:auc} to \N{gain:e3:lenient:auc} (+5.5 points) and its accuracy by 10.5 points.
That sits with Vela's middle results: request harm (+1.3 and +3.5 AUC points), hate speech (+12.3),
RAGTruth (+2.7). Vela's largest gains, +50 points of macro-F1 on fact-check-needed and feedback
routing, are against references that scored near chance (0.50 and 0.43), and its +19.4 on prompt
attacks is against a guard model on attack families it had not seen; there was much more room to
gain than Nox had on our tasks, where the released model already scored 79--84\%. On span tasks Vela
is mostly level with its references and loses 7.2 points on toxic spans, so not every target of a
broad fine-tune improves.

\paragraph{Accuracy is not the whole outcome.} On fact typing, the head tuned on other entities
raised accuracy by 4.7 points but left macro-F1 unchanged (\N{gain:e4facts:released:f1} to
\N{gain:e4facts:strict:f1}): the gain is on the common types, and the rare ones did not improve.
Even with the familiar entities in the training data, macro-F1 rose only 1.1 points. For a
cascade this is the useful kind of gain, because escalation is decided per item and common items
are most of the traffic, but it is narrower than the accuracy figure suggests, and it is the
metric Vela's card would have reported.

\paragraph{What the two cost.} Vela's card measures its own cost: on the Jev Decision Index of 38
general decision benchmarks it scores \N{vela:jev} against Nox's \N{vela:jev_nox}~\cite{vela2},
a quarter of its base's general-decision score. We did not measure the tuned heads outside our
tasks; they are not meant to serve any other question, and a tuned head changes every answer the
model gives, so in this pipeline it would be served only for these two tasks. The difference in
effort is the larger contrast: the head-only recipe needs cached features and labels the pipeline
already records, while Vela's needs a training run over the whole backbone.

\paragraph{Pre-tuned for another domain, Vela is a worse first stage here.} We also served the
three Vela decoder members, pinned and checked against their own recorded card outputs
(maximum difference below 0.001), on the frozen E3 and E4 items and replayed each as the cascade's
first stage under the same parity gate (Table~\ref{tab:velafirst}). None improves on the released
Nox it was tuned from. On triage every member must escalate more than half the tokens to keep real
names at qwen's level (Nox 33\%), and the cascade gains nothing; on fact typing Vela 4B escalates
less than Nox but its out-of-fold gate falls below qwen (\N{vf:e4facts:vela-4b:d_big} points, CI
\N{vf:e4facts:vela-4b:ci}), and the other members escalate most facts. The larger member is not
better: Vela 9B's confidence separates its right and wrong triage answers worst of all. Vela's
reference server is also slower than Nox's runtime (\N{vf:e3:vela-4b:small_ms} against
\N{vf:e3:nox:small_ms}~ms per triage item for the 4B), but replaying Vela's routing at Nox's
per-item latency leaves the conclusion unchanged on triage (1.48$\times$ against 2.07$\times$), and
on fact typing the accuracy shortfall remains whatever the speed. Removing Vela's built-in abstain
option, or escalating on its abstain probability, changed nothing material. Its tuning bought gains
on its own domain and cost it some of the general decision ability that our tasks draw on; tuning
toward our own rulings was the specialization that paid here.

\begin{table}[t]
\centering
\small
\setlength{\tabcolsep}{4pt}
\resizebox{\textwidth}{!}{\begin{tabular}{llrrrrrr}
\toprule
Task & First stage & Alone (\%) & Cascade (\%) & $\Delta$ vs qwen (pt) [95\% CI] & Esc.\ (\%) & ms/item & Speed-up \\
\midrule
E3 triage & Nox 4B, released & 79.2 & 84.7 & +0.0 [-0.3, +0.3] & 33 & 83 & 2.07$\times$ \\
 & Nox 4B, tuned & 89.7 & 87.5 & +2.8 [+1.7, +3.9] & 25 & 83 & 2.55$\times$ \\
 & Vela 2.0 0.8B & 75.6 & 84.7 & +0.0 [-0.7, +0.6] & 55 & 48 & 1.59$\times$ \\
 & Vela 2.0 4B & 76.7 & 84.6 & -0.1 [-0.3, +0.0] & 52 & 221 & 1.07$\times$ \\
 & Vela 2.0 9B & 77.8 & 84.5 & -0.2 [-0.7, +0.2] & 53 & 405 & 0.77$\times$ \\
\midrule
E4 fact typing & Nox 4B, released & 83.6 & 88.2 & -0.1 [-0.7, +0.5] & 40 & 94 & 1.70$\times$ \\
 & Nox 4B, tuned, disjoint & 88.4 & 88.9 & +0.6 [-0.6, +1.8] & 4 & 94 & 4.34$\times$ \\
 & Vela 2.0 0.8B & 68.2 & 87.9 & -0.4 [-0.7, -0.1] & 92 & 52 & 0.98$\times$ \\
 & Vela 2.0 4B & 83.6 & 87.4 & -0.9 [-1.7, -0.1] & 26 & 232 & 1.38$\times$ \\
 & Vela 2.0 9B & 82.5 & 88.1 & -0.2 [-0.6, +0.2] & 72 & 424 & 0.64$\times$ \\
\bottomrule
\end{tabular}
}
\caption{Vela 2.0 members as the cascade's first stage on our tasks, replay only (same parity gate,
2-fold by group, and qwen service times as Table~\ref{tab:replay}). ``Alone'': the first stage's own
accuracy. ms/item: the first stage's measured single-stream service time (Vela via its reference
server, Nox via its runtime). Tuned Nox rows repeated from Table~\ref{tab:replay} for comparison.
Vela served as released (abstain option on); removing it moved no figure by more than 0.4 points.}
\label{tab:velafirst}
\end{table}

\section{Discussion}

\paragraph{Specialization does shrink the second stage.} On triage, a head tuned on a few hundred of
the GM's rulings removes about a quarter of the escalations while holding the constraint that
matters, and on fact typing it closes the gap to the generative model for entities it has never
seen. For a bounded classification task that is the trade the question proposed, and it costs a
minute of CPU time rather than a GPU job, because the backbone is frozen and its features are cached.

\paragraph{But check what is being specialized to.} Fine-tuning on a pipeline's own review records
can learn two different things: the GM's \emph{conventions}, which transfer, and the GM's
\emph{rulings}, which are a table. E4 shows how much the second can masquerade as the first. A
leaderboard-style result (``tuned Nox removes 99\% of escalations at 5$\times$'') would have been
true on the replay and misleading as a claim about the model, and the fix is not a better model but
the companion study's conclusion: a ruling should be made once and read back. In this pipeline the
right fact-typing design is a lookup of the GM's past rulings first, tuned Nox for unseen entities,
and qwen only for the uncertain remainder.

\paragraph{The objective, not accuracy, decides.} Tuned Nox alone ``beats'' qwen on E3 by five points
of binary accuracy and still cannot replace it, because it drops more real names. Measuring the
first stage by escalation at the parity gate, with the expensive error constrained separately, is
what separated the useful variant from the two that looked better.

\paragraph{The replay extends to training.} The cascade study's replay chose among fixed first stages.
With cached features, the same machinery evaluates trained ones, as long as head and gate are
cross-fitted together. Seven variants were scored here in a few minutes of CPU time; only the two
that survived went to the hardware, and both did there what the replay said they would: the same
items escalated, the same accuracy difference within half a point, and the same speed-up within about
5\%. The replay's one blind spot was the one the cascade study named, contention on the second
stage, and it showed up as a failed measurement rather than a wrong prediction.

\section{Limitations}
\begin{itemize}
  \item \textbf{Head only.} The backbone was frozen. E6's LoRA predictions exist but use a different
        fold structure and were not cross-fitted with the gate; a LoRA first stage would also change
        Nox's latency and determinism, and was not tested.
  \item \textbf{Settings not tuned here, but not innocent.} Hyperparameters were fixed in advance from
        E6, which chose them on the same items with different folds.
  \item \textbf{Small label sets.} E3 has \N{e3:names} real names; the disjoint E3 failure is three
        names. E4's disjoint gate is fitted on few out-of-fold items.
  \item \textbf{Evaluated on the frozen items.} As in the cascade study, the live run uses the same
        items as the replay, with every label held out from its own head and gate; it is a test of the
        prediction, not of new sessions.
  \item \textbf{One GM, one campaign.} The specialization is to one person's conventions, which is
        the point for this pipeline and a limit for any wider claim.
  \item \textbf{Serving path.} The tuned head is swapped into a stock runtime after the released
        package verifies; the runtime's own integrity and golden checks therefore describe the released
        head, not the tuned one. Our per-head probe against the replay stands in for them.
\end{itemize}

\section{Conclusion}
Can a first stage tuned on the GM's rulings close the gap and remove the second stage? Partly, and
the replay says where before the hardware does. On fact typing, held out by entity, tuned Nox ties
the generative model, so the cascade sends it \N{live:e4facts:esc}\% of facts and runs
\N{live:e4facts:speedup_meas}$\times$ faster, but the larger gain on the GM's familiar entities is a
table, and is better kept as one. On triage the tuned model is more accurate than the generative
model and still not a replacement, because it drops more real names; as a first stage it cuts a
quarter of the escalations and makes the cascade both faster (\N{live:e3:speedup_meas}$\times$) and
more accurate than qwen alone. Specialization buys a cheaper cascade, not a single model, and the
measure that tells the two apart is the cascade's own objective, not the first stage's accuracy.
A published full fine-tune of the same base for another domain, Vela 2.0, gained about as much on
its mid-range targets as our head did on ours, and as a first stage for this pipeline it did worse
than the base it came from: what to specialize toward mattered more than how hard.

\section*{Data and code availability}
Head training (\texttt{clef/finetune/cascade\_head.py}; the evaluation scripts are in \texttt{clef/eval/}), the tuned-head server
(\texttt{serve\_tuned.py}), the cross-fitted replay and gate freezing
(\texttt{cascade\_ft.py}), the live driver (\texttt{cascade\_ft\_live.sh}), the
rerun of the contaminated legs (\texttt{cascade\_ft\_rerun.sh}), the Vela runs
(\texttt{vela\_run.sh}, \texttt{vela\_replay.py}, \texttt{serve\_vela.py}) and the script that generates every
number and table in this paper (\texttt{ft\_paper.py}) are in the public repository~\cite{dgxfun}. The campaign data is
not released because it quotes play sessions verbatim and names real people.

\section*{Use of AI assistance}
The experiments, analysis and a first draft of this paper were carried out by an AI coding agent
(Claude, Anthropic) working under the author's direction, with the author reviewing results and
making the decisions recorded in the paper. The author takes full responsibility for the content.

\bibliographystyle{plainnat}
\begin{thebibliography}{10}

\bibitem[Roussos(2026a)]{companion}
K.~Roussos.
\newblock Decision models at the game table: where non-generative classifiers help and fail in LLM
  tooling for tabletop role-playing campaigns.
\newblock Manuscript, October 2026. In~\cite{dgxfun}, \texttt{paper/main.tex}.

\bibitem[Roussos(2026b)]{cascade}
K.~Roussos.
\newblock Simulate before you serve: predicting a decision-model cascade from frozen outputs, and
  checking the prediction on real hardware.
\newblock Manuscript, October 2026. In~\cite{dgxfun}, \texttt{paper-cascade/main.tex}.

\bibitem[vLLM Semantic Router Team(2026)]{sysoneauto}
vLLM Semantic Router Team.
\newblock System One Auto.
\newblock Blog post, \url{https://vllm-sr.ai/blog/system-one-auto/}, accessed October 2026.

\bibitem[vLLM Semantic Router Team and KR Labs(2026)]{vela2}
vLLM Semantic Router Team and KR Labs.
\newblock Vela 2.0 (0.3B, 0.8B, 4B, 9B): model cards, evaluation and training notes.
\newblock \url{https://huggingface.co/vllm-sr/Vela-2.0-4B}, revision \texttt{d692394}, accessed October 2026.

\bibitem[Li et~al.(2018)]{l2sp}
X.~Li, Y.~Grandvalet, and F.~Davoine.
\newblock Explicit inductive bias for transfer learning with convolutional networks.
\newblock In \emph{ICML}, 2018.

\bibitem[Roussos(2026c)]{dgxfun}
K.~Roussos.
\newblock dgx-fun: experiments with local LLM serving on NVIDIA DGX Spark.
\newblock \url{https://github.com/kostadis/dgx-fun}, 2026.

\end{thebibliography}

\end{document}
